% Pursuing Stacks (cotes 134-2 to 134-8): archive cote 134-6, site page 106 (PDF page 107), typescript page 424.
% typescript. Draft compared: ps5.tex:7015-7092.
% AI-transcribed by Claude Opus 5.5 (claude-opus-5-5, adaptive thinking, effort high), Message Batches API, batch ps3 (msgbatch_01XdFmuVVBAcgc7EFxdtCqsj), 2026-11-12; not yet checked by eye.
\scanpage{134-6/106}
\mspage{424}

category \hand{$\mathcal{M}$} stable under small direct limits,   \unsure{\over{t}{T}}aking cosheaves on A\^{}
with values in \hand{$\mathcal{M}$},  and their restrictions to the six other categories in
the diagram (105),  we get a transposed seven term diagram as follows,
part of which reduces to the four terms diagram (58) (p.406) in case k=$\underline{Z}$ ;
\xout{xx}
\diagram{
\hand{(106)}\left\{\;
\begin{tikzcd}[sep=small]
\underline{\mathrm{Hom}}\hand{$_!$}(\mathrm{A}^{\wedge}\hand{$,\mathcal{M}$}) \ar[d,"\hand{$\approx$}"'] & & \\
\underline{\mathrm{Hom}}_{\hand{$k!$}}(\mathrm{A}^{\wedge}_{k},\hand{$\mathcal{M}$}) \ar[r,"\hand{$\approx$}"] &
\underline{\mathrm{Homaddinfkar}}_{k}(\mathrm{Proj}(\mathrm{A}^{\wedge}_{k}),\hand{$\mathcal{M}$}) \ar[r,dash,"\hand{$\approx$}"] \ar[d,"\hand{$\approx$}"] &
\underline{\mathrm{Homaddinf}}_{k}(\mathrm{Addinf}_{k}(\mathrm{A}),\hand{$\mathcal{M}$}) \ar[d,"\hand{$\approx$}"] \\
& \underline{\hand{H}\mathrm{omaddkar}}_{k}(\mathrm{UlProj}(\mathrm{A}^{\wedge}_{k}),\hand{$\mathcal{M}$}) \ar[r,dash,"\hand{$\approx$}"] &
\underline{\mathrm{Homadd}}_{k}(\mathrm{Add}_{k}(\mathrm{A}),\hand{$\mathcal{M}$}) \ar[d,"\hand{$\approx$}"] \\
& & \underline{\mathrm{Hom}}(\mathrm{A},\hand{$\mathcal{M}$}) \quad ,
\end{tikzcd}
\right.
}
% NOTE: the left brace of (106) is drawn by hand; a small mark is cut off at the right edge after Homadd_k(Add_k(A),M)

where the \xout{indexxk} meaning of the symbols used (such as index k,  suffixes
``add'' or ``addinf'' and ``kar'') for qualifying \uline{Hom}  and denoting various full
subcategories of \uline{Hom}  categories,  is clear from the explanations given
previously.  Replacing A by B and \uline{Hom}\hand{$_!$}(A\^{}, \hand{$\mathcal{M}$})  by \uline{Hom}\hand{$^!$}(\over{\unsure{A}}{B}\^{}, \hand{$\mathcal{M}$}),  we get a
diagram ``dual'' to (106)\hand{,} \hand{(}containing  the \xout{diagram} five-term diagram \hand{(}62)
(p.408)  in case k=$\underline{Z}$\hand{)},  wh\over{\unsure{ich}}{ich} we'll not write out here,  valid for any \inserted{k-additive}
category \hand{$\mathcal{M}$} stable under small inverse limits.   When \hand{$\mathcal{M}$} is a k-additive
category stable under both types of small limits,  then the last term of
the diagram (106)  is equal to the last term of the dual one,  hence
a system of fourteen mutually equivalent categ\over{\unsure{e}}{o}ries (compare p.408,  \unsure{where|when}
we considered ten among them only !),  expressing as many ways for \deleted{expres-}\inserted{interpre-}
\over{s}{t}ing the notion of an \hand{$\mathcal{M}$}-valued copresheaf on A i.e. an object of \uline{Hom}(A,\hand{$\mathcal{M}$})
(which is one among the fourteen ...).

Let's comme\over{\unsure{m}}{n}t a little on the significance of the various five \xout{xxxxx}
k-additive categories appearing in (105).  The largest one $\mathrm{A}^{\wedge}_{k}$  is there
precisely as the all-encompassing category of k-additive presheaves,
where to carry through all kinds of k-linear constructions between pre-%
sheaves on A.  The significance of the (second largest) subcategory
Pro\over{\unsure{f}}{j}($\mathrm{A}^{\wedge}_{k}$),  made up with all projective objects of $\mathrm{A}^{\wedge}_{k}$ ,  comes mainly from
homological algebra and emphasis upon replacing objects of $\mathrm{A}^{\wedge}_{k}$  by projec-%
tive resolutions;  \deleted{which}\inserted{these} are chain complexes in Proj($\mathrm{A}^{\wedge}_{k}$) ,  which may be
viewed as being defined (by any given object in $\mathrm{A}^{\wedge}_{k}$ ) ``up to chain homoto-%
py''.  M\over{\unsure{e}}{o}re sweepingly\deleted{,} still,  we get from general principles the canoni-%
cal equivalence of categories
\[ \deleted{\illegible{1}}\hand{(*)} \qquad \mathrm{D}^{-}(\mathrm{A}^{\wedge}_{k}) \hand{$\xleftarrow{\approx}$} \mathrm{K}^{-}(\mathrm{Proj}(\mathrm{A}^{\wedge}_{k})) \qquad , \]
where \deleted{\illegible{1}} D$^{-}$  designates the \deleted{\unsure{left}} ``derived category bounded from above''

\margin{106}
